\begin{proposition}[E.28 --- Proved: the pure boundary cubic component vanishes]
\label{E-no-pure-boundary-cubic}
Use the cubic notation of Propositions~\ref{E-cubic-symmetry}
and~\ref{E-five-parameters}. Rationally write the actual
Galois-stable splitting of homology as
\[
 V_{\mathbb Q_3}=H\oplus W,
 \qquad H=T_3E\otimes\mathbb Q_3,
 \qquad w=u\wedge v\in\bigwedge^2H.
\]
The image of the arithmetic class
$\eta\in H^1(G_{\mathbb Q,S_3},A_{\mathbb Q_3}(-1))$
under the canonical projection
\[
 A_{\mathbb Q_3}=\operatorname{Lie}_3(H\oplus W)
                  \longrightarrow\operatorname{Lie}_3(W)
 \tag{E.28.1}
\]
is zero. Consequently the projected symmetry-invariant class
of E.23 lies entirely in
\[
 H^1(G_{\mathbb Q,S_3},H(1)),
 \qquad \dim_{\mathbb Q_3}H^1(G_{\mathbb Q,S_3},H(1))=1.
 \tag{E.28.2}
\]
In the canonical rational direct sum of E.27, the actual
class therefore belongs to a four-dimensional subspace:
the three Kummer coordinates in
$H^1(G_{\mathbb Q,S_3},B^+_{\mathbb Q_3}(-1))$
and the single elliptic coordinate in (E.28.2).
The $\mathbb Q_3(2)\epsilon_K$ coordinate which was allowed
by symmetry alone is zero. These assertions determine a
vanishing projection, without assigning a value to the
remaining four coordinates or to integral torsion.
\end{proposition}
\begin{proof}
Let $\mathfrak n_{\leq3}$ be the rational class-three
Malcev Lie algebra of the free geometric pro-$3$ group
$N$. Its associated graded Lie algebra is the free Lie
algebra on $V_{\mathbb Q_3}$ through degree three.
The logarithm identifies
$(\gamma_2N/\gamma_4N)\otimes\mathbb Q_3$
with $\gamma_2\mathfrak n_{\leq3}$; this part is
abelian because a commutator of two degree-two terms
has degree at least four. This is an equivariant
identification, so it computes the extension defining
$\eta$ itself.

Choose lifts $a,b\in\mathfrak n_{\leq3}$ of $u,v$,
and choose lifts of a basis of $W$. They identify the
filtered free Lie algebra with its free graded model.
Set $\omega=[a,b]$. It is a lift of $w$. For
$g\in G_{\mathbb Q,S_3}$, stability of the actual
compact summand $H$ gives
\[
 ga=\alpha a+\beta b+q_a,
 \qquad gb=\gamma a+\delta b+q_b,
 \qquad q_a,q_b\in\gamma_2\mathfrak n_{\leq3},
\]
where $\alpha\delta-\beta\gamma=\chi_3(g)$.
Since $[q_a,q_b]=0$ in class three, the exact bracket
formula gives
\[
 g\omega-\chi_3(g)\omega
   =[\alpha a+\beta b,q_b]
       +[q_a,\gamma a+\delta b].
 \tag{E.28.3}
\]
Only the degree-two parts of $q_a,q_b$ contribute.
Every term of (E.28.3) therefore has at least one
factor in $H$. Under the multigrading it belongs to
\[
 \operatorname{Lie}_3(H)\oplus(H^2W)\oplus(HW^2),
\]
and its projection to $\operatorname{Lie}_3(W)$ is
zero. Dividing (E.28.3) by $\chi_3(g)$ gives the
cocycle of the extension defining $\eta$. A different
lift of $w$ changes that cocycle by a coboundary in
$A_{\mathbb Q_3}(-1)$. Thus (E.28.1) kills its
cohomology class independently of all these lifts.

The submodule $B=[V,w]$ also has at least two compact
factors in each bracket, so (E.28.1) factors through
$Q=A/B$. The geometric group
$\mathcal G=E[3]\rtimes\mu_6$ preserves the actual
homology splitting $H\oplus W$. Its rational
averaging projector therefore commutes with the
multigrading projections. E.23 identifies
\[
 (Q_{\mathbb Q_3})^{\mathcal G}
       =D_{HW^2}\oplus D_{W^3},
 \qquad
 D_{HW^2}\simeq H(2),
 \quad
 D_{W^3}\simeq\mathbb Q_3(3)
                     \oplus\mathbb Q_3(3)\epsilon_K.
\]
These are Galois-stable direct summands. The zero
projection just proved kills the whole $D_{W^3}$
component of the projected class. Its remaining
component, after twist $(-1)$, is exactly (E.28.2).
The cohomology dimension one is the same Euler
characteristic and $H^2$-vanishing computation used
for this summand in E.23.

Finally $B^{\mathcal G}_{\mathbb Q_3}=0$, so the
averaging projector on $A_{\mathbb Q_3}$ identifies
$A^{\mathcal G}_{\mathbb Q_3}$ with
$Q^{\mathcal G}_{\mathbb Q_3}$, as established in
E.27. That projector also preserves the multigrading.
Its canonical direct-sum description therefore
restricts to the three-dimensional Kummer kernel
and the one-dimensional summand (E.28.2), giving
dimension four. This uses the rational compact
splitting and rational Malcev logarithm; it imposes
no unproved integral vanishing on the torsion
recorded in E.24 and E.27.
\end{proof}
